\clearpage
\SourceLeaf{147}
\begin{center}
\textsc{Gauss} to \textsc{Olbers}. Göttingen, 24 January 1812.
\end{center}

\noindent . . . . . . . . . . . I hear that Mr.\ \textsc{Delambre} has written a very long-winded rhetorical exercise, in his own fashion, on the moindres carrés in the French\ Moniteur. I do not read the Moniteur, which in any case reaches us only in annual installments. Perhaps you will have an opportunity sometime to attest \textls{publicly} that I explained its main points to you as early as our first meeting in person in 1803. Among my papers I find that in June 1798, when I had \textls{long} been using that method, I first saw \textsc{Laplace}'s method and noted its incompatibility with the principles of the calculus of probabilities in a brief journal\HistoricalNote{[*)]}{[*)}{In the diary of \textsc{Gauss} already mentioned on page 20, the following entry appears for 1798: »\emph{Calculus probabilitatis contra La Place defensus. Gott. Jun. 17}«.]} of my mathematical pursuits. In the autumn of 1802 I entered the \textsc{viii}\textsuperscript{th} set of elements of Ceres in my astronomical draft notebook, having obtained them by the method of least squares. The papers in which I applied that method in earlier years, e.\,g.\ in the spring of 1799, to \textsc{Ulugh Beigh}'s table of the equation of time have been lost. The only wonder is that this principle, which suggests itself so readily that one can set no special value on the idea alone, was not applied 50 or 100 years earlier by others, e.\,g.\ by \textsc{Euler}, \textsc{Lambert}, \textsc{Halley}, or \textsc{Tobias Mayer}, though it could very well be that, e.\,g.\ the last of these applied something of the kind without proclaiming it, just as every reckoner necessarily devises for himself many shortcuts and methods that he passes on only occasionally by word of mouth.

\noindent . . . . . . . . . . . . .

